
\begin{figure}[htbp]
    \centering

    \begin{tikzpicture}[
        node distance=1.5cm and 1.8cm,
        every node/.style={font=\small},
        central/.style={
            rectangle,
            rounded corners=4pt,
            draw=black,
            thick,
            minimum width=5.2cm,
            minimum height=1.15cm,
            align=center,
            fill=blue!8
        },
        dimension/.style={
            rectangle,
            rounded corners=3pt,
            draw=black,
            minimum width=3.5cm,
            minimum height=0.9cm,
            align=center,
            fill=gray!10
        },
        distinction/.style={
            rectangle,
            rounded corners=3pt,
            draw=black,
            minimum width=3.5cm,
            minimum height=0.9cm,
            align=center,
            fill=orange!10
        },
        relation/.style={
            <->,
            thick,
            >=stealth
        },
        derivation/.style={
            ->,
            thick,
            >=stealth
        }
    ]

    % Central ontological unity
    \node[central] (beingidentity) {
        \textbf{ONTOLOGICAL ACTUALITY}\\
        Being and Identity
    };

    % Analytical dimensions
    \node[dimension, below left=of beingidentity] (being) {
        \textbf{Being}\\
        The actuality of an entity
    };

    \node[dimension, below right=of beingidentity] (identity) {
        \textbf{Identity}\\
        The entity being itself
    };

    % Constitutive determination
    \node[central, below=2.0cm of beingidentity] (determination) {
        \textbf{DETERMINATE CONSTITUTION}\\
        Constitutive determinations of the entity
    };

    % Analytical relations
    \node[distinction, below left=1.5cm and 0.5cm of determination] (unity) {
        \textbf{Unity}\\
        Ontological coherence
    };

    \node[distinction, below=1.5cm of determination] (individuation) {
        \textbf{Individuation}\\
        Differentiation as this entity
    };

    \node[distinction, below right=1.5cm and 0.5cm of determination] (distinction) {
        \textbf{Distinction}\\
        Difference from others
    };

    % Non-equivalent characterizations
    \node[dimension, below=1.8cm of individuation] (signature) {
        \textbf{Signature}\\
        Structured characterization
    };

    \node[dimension, left=of signature] (invariance) {
        \textbf{Invariance}\\
        Persistence across transformations
    };

    \node[dimension, right=of signature] (manifestation) {
        \textbf{Manifestation}\\
        Ontological realization
    };

    % Relations between being and identity
    \draw[relation] (beingidentity.south west) -- (being.north);
    \draw[relation] (beingidentity.south east) -- (identity.north);

    \draw[relation] (being) -- (determination);
    \draw[relation] (identity) -- (determination);

    % Constitutive and analytical dimensions
    \draw[derivation] (determination.south west) -- (unity.north);
    \draw[derivation] (determination.south) -- (individuation.north);
    \draw[derivation] (determination.south east) -- (distinction.north);

    % Characterizations and related concepts
    \draw[relation] (individuation.south west) -- (invariance.north);
    \draw[relation] (individuation.south) -- (signature.north);
    \draw[relation] (individuation.south east) -- (manifestation.north);

    \end{tikzpicture}

    \caption{The Ontological Structure of Identity: The Inseparability of Being and Identity and Their Analytical Relations to Determination, Unity, Individuation, and Related Concepts}
    \label{fig:identity-ontological-structure}
\end{figure}
